\chapter{Estado del arte}\label{chap:estado_del_arte}

Este capítulo revisa los trabajos relacionados con la aplicación de la inteligencia artificial a la imagen médica y con la predicción de complicaciones tras la biopsia pulmonar. También se examinan los avances en radiómica, \textit{deep learning}, \textit{PU learning} y descubrimiento de subgrupos que permiten situar la propuesta desarrollada en este trabajo.

\section{Inteligencia artificial en medicina e imagen médica}

La inteligencia artificial se ha consolidado en los últimos años como una herramienta de apoyo a la práctica médica, especialmente en tareas que implican el análisis de grandes volúmenes de datos clínicos, imágenes médicas y registros electrónicos. Su interés en medicina no reside únicamente en automatizar tareas diagnósticas, sino también en mejorar la estratificación de riesgo, apoyar la toma de decisiones y avanzar hacia una medicina más personalizada. En este contexto, diversos trabajos de revisión han señalado que la incorporación de sistemas de IA plantea oportunidades relevantes, pero también desafíos relacionados con la calidad de los datos, la validación clínica, la interpretabilidad, la integración en el flujo asistencial y las implicaciones éticas y legales \cite{Yu2018,Topol2019}.

Dentro de las áreas médicas, la imagen médica ha sido uno de los campos donde la IA ha mostrado un desarrollo más rápido. Modalidades como la radiografía, la tomografía computarizada, la resonancia magnética, la ecografía o la histopatología digital generan datos visuales de alta dimensionalidad, cuya interpretación depende de patrones espaciales complejos. Las técnicas de \textit{deep learning}, especialmente las redes neuronales convolucionales, han permitido pasar de enfoques basados en características diseñadas manualmente a modelos capaces de aprender representaciones directamente a partir de los datos. Litjens et al. \cite{Litjens2017} revisaron más de 300 contribuciones en imagen médica y mostraron la aplicación de \textit{deep learning} en tareas como clasificación, detección, segmentación, registro y reconstrucción de imágenes.

Uno de los trabajos que contribuyó a visibilizar el potencial clínico de \textit{deep learning} fue el de Esteva et al. \cite{Esteva2017}, donde una red neuronal convolucional entrenada con 129.450 imágenes clínicas de lesiones cutáneas alcanzó un rendimiento comparable al de dermatólogos expertos en tareas binarias de clasificación de cáncer de piel. Aunque este estudio pertenece al ámbito dermatológico, resulta relevante como ejemplo paradigmático de un modelo entrenado extremo a extremo a partir de imágenes y etiquetas clínicas, sin necesidad de definir manualmente las características visuales relevantes.

En radiología, y especialmente en tomografía computarizada, la IA permite analizar información tridimensional y patrones espaciales complejos que pueden ser difíciles de sintetizar mediante inspección visual convencional. Esta capacidad ha impulsado el desarrollo de modelos orientados a tareas como detección de lesiones, segmentación anatómica, clasificación diagnóstica y estratificación de riesgo, consolidando la imagen médica como uno de los ámbitos de mayor aplicación de \textit{deep learning}.

Una línea reciente de investigación se orienta hacia modelos fundacionales o generalistas, entrenados sobre grandes conjuntos de datos y adaptables a múltiples tareas clínicas con poca supervisión adicional. Moor et al. \cite{Moor2023} plantean el concepto de \textit{Generalist Medical AI}, referido a modelos capaces de integrar distintas modalidades médicas, como imagen, texto clínico, registros electrónicos, resultados de laboratorio o genómica, y de generar salidas flexibles como explicaciones, anotaciones o recomendaciones. En imagen médica, esta tendencia también se observa en modelos de segmentación generalista como MedSAM \cite{Ma2023}, entrenado con más de 1,5 millones de pares imagen-máscara procedentes de múltiples modalidades, o herramientas como TotalSegmentator \cite{wasserthal2023totalsegmentator}, orientadas a la segmentación automática robusta de estructuras anatómicas en TC. Estos enfoques resultan particularmente útiles en contextos clínicos con pocos datos etiquetados, ya que permiten aprovechar representaciones preentrenadas, segmentaciones anatómicas y estrategias de transferencia para abordar tareas predictivas específicas.

En conjunto, estos avances muestran que la inteligencia artificial ha pasado de ser una herramienta centrada en tareas diagnósticas concretas a convertirse en un marco metodológico capaz de integrar información heterogénea, extraer patrones complejos y apoyar procesos de decisión clínica. En imagen médica, esta evolución amplía las posibilidades de combinar modelos predictivos, segmentación automática y estrategias de transferencia para abordar problemas clínicos con datos limitados y alta complejidad.

\section{Biopsia pulmonar guiada por TC: contexto clínico y predicción de complicaciones}

Como se comentó previamente, el cáncer de pulmón constituye uno de los principales problemas de salud pública a nivel mundial, tanto por su elevada incidencia como por su alta mortalidad. Según las estimaciones de GLOBOCAN 2022 \cite{Bray2024}, en el año 2022 se registraron aproximadamente 20 millones de nuevos casos de cáncer y 9,7 millones de muertes por cáncer en todo el mundo. Dentro de este escenario, el cáncer de pulmón fue el tumor diagnosticado con mayor frecuencia, con alrededor de 2,5 millones de nuevos casos, lo que representa aproximadamente el 12,4\% de todos los diagnósticos oncológicos. Además, fue la principal causa de muerte por cáncer, con cerca de 1,8 millones de fallecimientos, equivalentes al 18,7\% de la mortalidad oncológica global. En este contexto, la TC torácica desempeña un papel fundamental en la detección, caracterización y seguimiento de lesiones pulmonares, así como en la planificación del proceso diagnóstico.

Como se ha expuesto en el capítulo anterior, la biopsia pulmonar guiada por TC combina un elevado rendimiento diagnóstico con un riesgo de complicaciones que debe considerarse durante la planificación del procedimiento. Su rendimiento diagnóstico se ha situado en torno al 90-93\%. No obstante, tanto la capacidad diagnóstica como la frecuencia de eventos adversos pueden variar en función de la técnica empleada. En el estudio de Ho et al. \cite{Ho2023}, la biopsia transtorácica guiada por TC alcanzó un rendimiento diagnóstico del 83,45\%, frente al 68,82\% de la biopsia transbronquial guiada por ecografía endobronquial radial. Esta diferencia estuvo acompañada de mayores tasas de neumotórax, 21,43\% frente a 2,87\%, y de hemorragia, 9,0\% frente a 5,7\%. Por otra parte, el metaanálisis de Heerink et al. \cite{Heerink2017} estimó tasas globales de complicaciones de aproximadamente el 39\% para las biopsias con aguja gruesa y del 24\% para la aspiración con aguja fina, e identificó el neumotórax y la hemorragia pulmonar como los eventos adversos más frecuentes.

Más allá de la técnica empleada, la aparición de complicaciones está condicionada por factores relacionados con el paciente, la lesión y el propio procedimiento. Entre ellos se encuentran la presencia de enfisema o enfermedad pulmonar obstructiva crónica (EPOC), el tamaño y la profundidad de la lesión, su contacto con la pleura, la posición del paciente y el trayecto de la aguja. Esta diversidad de factores explica que el riesgo no sea uniforme entre pacientes y justifica el desarrollo de herramientas capaces de estimarlo antes de la intervención. Una predicción individualizada podría contribuir a identificar a los pacientes con mayor riesgo, apoyar la planificación del procedimiento y favorecer una toma de decisiones clínicas más personalizada.

La predicción de complicaciones asociadas a la biopsia pulmonar ha sido abordada en la literatura principalmente desde enfoques estadísticos basados en factores clínicos, anatómicos y procedimentales. En particular, el neumotórax ha recibido mayor atención que otras complicaciones debido a su elevada frecuencia tras la biopsia. Algunos trabajos han propuesto modelos de riesgo o nomogramas para estimar la probabilidad de neumotórax tras biopsia guiada por TC \cite{Anzidei2015,Zhao2017,Wang2019,Weon2021,Yang2021}. Además, Zhao et al. \cite{Zhao2022} desarrollaron y validaron un modelo basado en regresión logística para predecir neumotórax tras biopsia pulmonar coaxial guiada por TC, incorporando ocho factores de riesgo: sexo, posición del paciente, campo pulmonar, contacto de la lesión con la pleura, tamaño de la lesión, distancia pleura-lesión, presencia de enfisema adyacente al trayecto de biopsia y cruce de cisuras. El modelo mostró una capacidad predictiva moderada, con AUC de 0,71 en entrenamiento y validación, y 0,68 en el conjunto de test interno.

Más recientemente, Yang et al. \cite{yang2026machine} desarrollaron un sistema dual basado en aprendizaje automático para pacientes sometidos a biopsia pulmonar percutánea guiada por TC, incluyendo un segundo análisis específico para aquellos con una biopsia transbronquial previa no diagnóstica. A diferencia de los métodos que procesan directamente las imágenes, los modelos utilizaron exclusivamente variables tabulares. Aunque las TC se emplearon para guiar el procedimiento y obtener información sobre la lesión y el trayecto de la aguja, las imágenes no se introdujeron directamente en los algoritmos. En su lugar, la información radiológica se incorporó mediante descriptores medidos y registrados por el equipo médico, como el contacto pleural, el tamaño y las características de la lesión o la profundidad de punción, junto con variables clínicas y procedimentales. El estudio multicéntrico incluyó 2910 biopsias percutáneas para predecir eventos adversos graves, definidos principalmente como neumotórax con necesidad de drenaje o hemorragia de alto grado. Por otra parte, sobre el subconjunto de pacientes con una biopsia transbronquial previa no diagnóstica se estimó la probabilidad de éxito diagnóstico de la biopsia percutánea de rescate. Tras seleccionar las variables mediante LASSO, los autores compararon diez algoritmos de aprendizaje automático, entre ellos regresión logística, SVM, Random Forest, redes neuronales, LightGBM, CatBoost y XGBoost. El modelo basado en XGBoost obtuvo una AUC externa de 0,717 para la predicción de eventos adversos graves y de 0,884 para la predicción del éxito diagnóstico. Finalmente, el análisis mediante SHAP permitió identificar la contribución de variables como el contacto pleural, la profundidad de punción, el tamaño de la lesión, el sexo y el tabaquismo en la predicción de los eventos adversos graves.

Estos antecedentes ponen de manifiesto que la estimación preprocedimiento del riesgo de complicaciones es una línea de investigación activa y clínicamente relevante. Aun así, gran parte de los estudios existentes se centra en una única complicación, especialmente el neumotórax, o en eventos adversos graves dentro de escenarios clínicos concretos. Otras complicaciones, como la hemorragia pulmonar o la hemoptisis, han sido estudiadas principalmente mediante análisis de factores de riesgo, con menor desarrollo de modelos predictivos específicos \cite{Yeow2004,Zhu2020}. Esta situación motiva la exploración de enfoques más amplios que integren variables clínicas e información cuantitativa derivada de la TC para predecir la aparición de complicaciones y caracterizar sus distintos tipos.

\section{Aprendizaje automático y radiómica en imagen médica}

La radiómica surge como una aproximación orientada a transformar las imágenes médicas en datos cuantitativos explotables mediante técnicas estadísticas y de aprendizaje automático. A diferencia de la interpretación visual convencional, este enfoque extrae un conjunto amplio de características que describen propiedades de intensidad, forma y textura de regiones de interés, con el objetivo de capturar patrones no siempre perceptibles por el observador humano. Lambin et al. \cite{Lambin2012} introdujeron formalmente este paradigma como una estrategia para extraer más información de las imágenes médicas mediante análisis avanzado de características, mientras que Gillies et al. \cite{Gillies2016} popularizaron la idea de que las imágenes médicas no son únicamente representaciones visuales, sino también fuentes de datos cuantificables.

Uno de los trabajos de referencia en este campo fue el de Aerts et al. \cite{Aerts2014}, quienes mostraron que las características radiómicas extraídas de TC podían presentar valor pronóstico en cohortes independientes de pacientes con cáncer de pulmón y con tumores de cabeza y cuello. Desde entonces, la radiómica se ha aplicado en múltiples tareas de imagen médica, incluyendo diagnóstico, pronóstico, estratificación de pacientes y predicción de respuesta terapéutica.

Pese a su potencial, la radiómica presenta importantes desafíos metodológicos. La extracción de características puede verse afectada por la modalidad de imagen, el protocolo de adquisición, la reconstrucción, el preprocesamiento, la segmentación y la discretización de intensidades. Por ello, la reproducibilidad se ha convertido en uno de los principales retos del área. Iniciativas como la Image Biomarker Standardisation Initiative (IBSI) han propuesto definiciones estandarizadas y valores de referencia para características radiómicas, con el objetivo de mejorar la comparabilidad entre estudios y herramientas \cite{Zwanenburg2020}. Del mismo modo, herramientas abiertas como PyRadiomics han contribuido a sistematizar la extracción de características y facilitar la reproducibilidad de los experimentos \cite{VanGriethuysen2017}.

En los últimos años, la radiómica ha evolucionado hacia enfoques híbridos que combinan características manuales, aprendizaje automático clásico, \textit{deep learning} y estrategias de selección automática de modelos. Revisiones recientes destacan que los métodos de aprendizaje automático y \textit{deep learning} pueden mejorar la explotación de las características radiómicas, aunque insisten en la necesidad de validación externa, control de fuga de información y análisis de estabilidad de las variables \cite{Avanzo2020,Mayerhoefer2020}. Esta consideración es especialmente relevante en problemas con pocos pacientes y alta dimensionalidad, donde los modelos clásicos basados en características radiómicas pueden ofrecer una alternativa más robusta e interpretable que modelos profundos entrenados.

Por estas razones, la radiómica resulta especialmente útil en escenarios clínicos con cohortes reducidas y datos de imagen de alta dimensionalidad. Al transformar la TC en un conjunto estructurado de descriptores cuantitativos, permite combinar información derivada de la imagen con variables clínicas y aplicar modelos de aprendizaje automático relativamente interpretables. Esta combinación la convierte en una alternativa sólida frente a enfoques profundos cuando el tamaño muestral es limitado y la explicabilidad del modelo es un requisito importante.

\section{\textit{Deep learning} en TC torácica}

El \textit{deep learning} ha adquirido un papel destacado en el análisis de TC torácica, especialmente en tareas relacionadas con la detección, segmentación y caracterización de nódulos pulmonares. A diferencia de los enfoques radiómicos clásicos, que dependen de características definidas manualmente, los modelos profundos aprenden representaciones directamente a partir de los datos de imagen. Esto resulta especialmente relevante en TC, donde la información es volumétrica y los patrones clínicamente significativos pueden depender tanto de la apariencia local de la lesión como del contexto anatómico circundante.

Uno de los ámbitos donde el \textit{deep learning} ha mostrado mayor desarrollo es la detección automática de nódulos pulmonares. Retos como LUNA16 \cite{Setio2017} impulsaron la comparación sistemática de algoritmos sobre bases de datos públicas de TC torácica, mostrando el potencial de las redes convolucionales para identificar candidatos nodulares y reducir falsos positivos. Estos trabajos consolidaron el uso de arquitecturas profundas como herramientas de apoyo al cribado y al diagnóstico radiológico, especialmente en escenarios con un gran volumen de imágenes y una interpretación radiológica compleja.

Una línea especialmente relevante es el uso de modelos tridimensionales extremo a extremo para cribado de cáncer de pulmón. Ardila et al. \cite{Ardila2019} propusieron un modelo de \textit{deep learning} 3D capaz de analizar volúmenes completos de TC de baja dosis y estimar el riesgo de cáncer de pulmón, alcanzando un AUC de 0,944 en casos del National Lung Cancer Screening Trial y mostrando un rendimiento comparable al de radiólogos expertos en determinados escenarios. Más recientemente, modelos como Sybil han extendido esta idea hacia la predicción longitudinal del riesgo futuro de cáncer de pulmón a partir de una única TC de baja dosis, reforzando el interés por utilizar la imagen torácica no solo para diagnóstico inmediato, sino también para estratificación de riesgo \cite{Mikhael2023}.

El desarrollo de modelos profundos en TC torácica también plantea dificultades metodológicas relevantes. Su rendimiento depende en gran medida del tamaño y diversidad de los datos disponibles, de la calidad de las anotaciones y de la homogeneidad de los protocolos de adquisición. Diferencias entre hospitales, escáneres, reconstrucciones, tamaños de vóxel o criterios de inclusión pueden afectar de forma notable a la generalización. Además, la interpretación de las representaciones aprendidas sigue siendo compleja, por lo que técnicas como Grad-CAM, mapas de atención o métodos basados en atribución resultan necesarias para valorar si el modelo se apoya en regiones clínicamente coherentes.

Estas limitaciones adquieren especial relevancia en problemas menos estudiados que la detección o clasificación de nódulos pulmonares. En el caso de la predicción de complicaciones tras una biopsia pulmonar, el evento adverso no depende únicamente de la apariencia de la lesión en la TC, sino también de factores anatómicos, clínicos y procedimentales. Por ello, los modelos profundos deben evaluarse de forma comparativa frente a enfoques radiómicos y modelos clásicos, prestando especial atención al tamaño muestral, la integración de variables clínicas, la interpretabilidad y la validación externa.

\section{\textit{PU learning} en contextos clínicos}

En los problemas clínicos reales, las etiquetas disponibles no siempre representan de forma completa o inequívoca el estado del paciente. La ausencia de una etiqueta diagnóstica o de un evento adverso registrado no implica necesariamente que dicho evento no haya ocurrido, sino que puede deberse a limitaciones en el registro, seguimiento incompleto o criterios de anotación heterogéneos. Este escenario se relaciona con el aprendizaje débilmente supervisado y, en particular, con el \textit{PU learning}, donde se dispone de ejemplos positivos confirmados y de un conjunto no etiquetado que puede contener tanto positivos no identificados como verdaderos negativos. Elkan y Noto formalizaron este planteamiento para entrenar clasificadores a partir de datos positivos y no etiquetados \cite{Elkan2008}, y revisiones posteriores han sistematizado sus principales supuestos, estrategias de estimación y ámbitos de aplicación \cite{Bekker2020}.

El interés del \textit{PU learning} en medicina se debe a que muchos conjuntos clínicos presentan anotaciones incompletas o sesgadas hacia los casos más evidentes. En este contexto, considerar todos los ejemplos no etiquetados como negativos puede introducir ruido y afectar a la estimación del riesgo. Los métodos de \textit{PU learning} intentan abordar esta limitación mediante distintas estrategias, como la identificación de negativos fiables, la estimación de la probabilidad a priori de la clase positiva o la formulación de funciones de riesgo adaptadas. En particular, Kiryo et al. \cite{Kiryo2017} propusieron un estimador de riesgo no negativo para \textit{PU learning}, diseñado para reducir problemas de sobreajuste cuando se emplean modelos flexibles, como redes neuronales. Su aplicación, sin embargo, depende de los supuestos adoptados sobre el mecanismo de etiquetado y sobre la representatividad de los positivos disponibles.

\section{Descubrimiento de subgrupos en medicina}

Además de predecir el riesgo individual, en medicina resulta especialmente relevante identificar perfiles de pacientes que compartan patrones clínicos, anatómicos o procedimentales asociados a un comportamiento diferencial. En este sentido, el descubrimiento de subgrupos constituye una técnica de minería de datos orientada a encontrar subgrupos interpretables respecto a una variable objetivo. A diferencia de los modelos predictivos globales, cuyo objetivo principal es optimizar el rendimiento sobre toda la población, el descubrimiento de subgrupos busca reglas descriptivas del tipo condición--resultado, priorizando su interpretabilidad, cobertura y relevancia estadística. Este enfoque ha sido ampliamente estudiado como técnica de análisis exploratorio supervisado \cite{Herrera2011,Atzmueller2015} y ha mostrado interés en contextos biomédicos para la identificación de perfiles de riesgo y la generación de hipótesis clínicas \cite{Gamberger2002}.

En biopsia pulmonar, este enfoque puede aportar una lectura complementaria a los modelos predictivos, al permitir explorar combinaciones de variables asociadas a una mayor frecuencia de complicaciones. Su utilidad dependerá de que los subgrupos encontrados sean suficientemente estables, clínicamente plausibles y no derivados únicamente del azar.

% \section{Explicabilidad, validación y reproducibilidad en IA médica???}

% \section{Síntesis del estado del arte}
